% GENERATED FILE. Edit canonical lessons, then run npm run generate:books.
% canonical-source-hash: fnv1a64:925ef3b01ba1feda
% canonical-lessons: BN-C447-aphsos, BN-C447-onushochona, BN-C447-onutopto, BN-C447-postano, BN-C447-svikar-kora

\chapter{Regret}
\label{ch:bn-regret}
\hlchaptermodality{447}

\begin{chapteropening}
\textbf{By the end of this chapter:} \emph{I can say what I regret and admit a mistake.}

Complete the last lesson of chapter 447: I can say what I regret and admit a mistake.
\end{chapteropening}

\section[\texorpdfstring{\bn{আফসোস} (āphsos)}{আফসোস (āphsos)}]{\bn{আফসোস} (āphsos) — regret; a pity}

\label{lesson:BN-C447-aphsos}



\begin{quote}
Before the new one: say the Bengali for to breathe a sigh of relief, then the Bengali for to cope with.
\end{quote}

\subsection*{You'll want to know: \bn{আফসোস}}
\textbf{\bn{আফসোস}} — \emph{āphsos} — \textquotedblleft{}regret; a pity\textquotedblright{}.

\begin{quote}
\bn{ছুটিতে} \bn{গ্রামে} \bn{যেতে} \bn{পারলাম} \bn{না}, \bn{এটাই} \bn{আমার} \bn{আফসোস।}
\end{quote}

\begin{quote}
\emph{chhuṭite grāme jete pārlām nā, eṭāi āmār āphsos.}
\end{quote}

\begin{quote}
I couldn't go to the village for the holidays, and that is my one regret.
\end{quote}

\bn{আফসোস} \bn{করা}, to say you regret something; \bn{আফসোস} \bn{হওয়া}, to feel regret.

\begin{cousinweb}
\textbf{\bn{আফসোস}} (\emph{āphsos}) is usually traced to Persian \emph{afsōs}, a cry of \textquotedblleft{}alas\textquotedblright{} that also means regret.
\end{cousinweb}

\begin{grammarlens}[title={What you've built}]
One more word for describing what you lived and felt.
\end{grammarlens}

\subsection*{Your turn}
\begin{itemize}
\raggedright
  \item \emph{Say it:} \emph{āphsos}
  \item \emph{Say it:} \emph{āphsos}, once more
  \item \emph{Say it:} \emph{chhuṭite grāme jete pārlām nā, eṭāi āmār āphsos.} — I couldn't go to the village for the holidays, and that is my one regret.
  \item \emph{Recall:} say the Bengali for to breathe a sigh of relief, then the Bengali for to cope with, then say \emph{āphsos} again
\end{itemize}

\begin{tcolorbox}[breakable,colback=teal!4,colframe=teal!35!black,title={Before you move on}]
What does \bn{আফসোস} mean? (\textquotedblleft{}Regret; a pity\textquotedblright{}.) Say it once more.
\end{tcolorbox}
\section[\texorpdfstring{\bn{অনুশোচনা} (ônushochônā)}{অনুশোচনা (ônushochônā)}]{\bn{অনুশোচনা} (ônushochônā) — remorse; repentance}

\label{lesson:BN-C447-onushochona}



\begin{quote}
Before the new one: say the Bengali for to cope with, then the Bengali for regret.
\end{quote}

\subsection*{You'll want to know: \bn{অনুশোচনা}}
\textbf{\bn{অনুশোচনা}} — \emph{ônushochônā} — \textquotedblleft{}remorse; repentance\textquotedblright{}.

\begin{quote}
\bn{মিথ্যে} \bn{বলার} \bn{পরে} \bn{অরুণের} \bn{অনুশোচনা} \bn{হল।}
\end{quote}

\begin{quote}
\emph{mitthe bôlār pôre ôruner ônushochônā hôlô.}
\end{quote}

\begin{quote}
After telling the lie, Arun felt remorse.
\end{quote}

\begin{cousinweb}
\textbf{\bn{অনুশোচনা}} (\emph{ônushochônā}) is a tatsama, Sanskrit \emph{anuśocana}, \textquotedblleft{}sorrow, repentance\textquotedblright{} (Monier-Williams), grief that comes after the deed.
\end{cousinweb}

\begin{grammarlens}[title={What you've built}]
One more word for describing what you lived and felt.
\end{grammarlens}

\subsection*{Your turn}
\begin{itemize}
\raggedright
  \item \emph{Say it:} \emph{ônushochônā}
  \item \emph{Say it:} \emph{ônushochônā}, once more
  \item \emph{Say it:} \emph{mitthe bôlār pôre ôruner ônushochônā hôlô.} — After telling the lie, Arun felt remorse.
  \item \emph{Recall:} say the Bengali for to cope with, then the Bengali for regret, then say \emph{ônushochônā} again
\end{itemize}

\begin{tcolorbox}[breakable,colback=teal!4,colframe=teal!35!black,title={Before you move on}]
What does \bn{অনুশোচনা} mean? (\textquotedblleft{}Remorse; repentance\textquotedblright{}.) Say it once more.
\end{tcolorbox}
\section[\texorpdfstring{\bn{অনুতপ্ত} (ônutôptô)}{অনুতপ্ত (ônutôptô)}]{\bn{অনুতপ্ত} (ônutôptô) — repentant; contrite}

\label{lesson:BN-C447-onutopto}



\begin{quote}
Before the new one: say the Bengali for regret, then the Bengali for remorse.
\end{quote}

\subsection*{You'll want to know: \bn{অনুতপ্ত}}
\textbf{\bn{অনুতপ্ত}} — \emph{ônutôptô} — \textquotedblleft{}repentant; contrite\textquotedblright{}.

\begin{quote}
\bn{চোরটা} \bn{অনুতপ্ত} \bn{হয়ে} \bn{টাকা} \bn{ফিরিয়ে} \bn{দিল।}
\end{quote}

\begin{quote}
\emph{chorṭā ônutôptô hôye ṭākā phiriye dilô.}
\end{quote}

\begin{quote}
The thief, repentant, gave the money back.
\end{quote}

\begin{cousinweb}
\textbf{\bn{অনুতপ্ত}} (\emph{ônutôptô}) is a tatsama, Sanskrit \emph{anutapta}, \textquotedblleft{}heated\textquotedblright{} and so \textquotedblleft{}filled with regret\textquotedblright{} (Monier-Williams): burning afterwards.
\end{cousinweb}

\begin{grammarlens}[title={What you've built}]
One more word for describing what you lived and felt.
\end{grammarlens}

\subsection*{Your turn}
\begin{itemize}
\raggedright
  \item \emph{Say it:} \emph{ônutôptô}
  \item \emph{Say it:} \emph{ônutôptô}, once more
  \item \emph{Say it:} \emph{chorṭā ônutôptô hôye ṭākā phiriye dilô.} — The thief, repentant, gave the money back.
  \item \emph{Recall:} say the Bengali for regret, then the Bengali for remorse, then say \emph{ônutôptô} again
\end{itemize}

\begin{tcolorbox}[breakable,colback=teal!4,colframe=teal!35!black,title={Before you move on}]
What does \bn{অনুতপ্ত} mean? (\textquotedblleft{}Repentant; contrite\textquotedblright{}.) Say it once more.
\end{tcolorbox}
\section[\texorpdfstring{\bn{পস্তানো} (pôstāno)}{পস্তানো (pôstāno)}]{\bn{পস্তানো} (pôstāno) — to rue; to regret (something done)}

\label{lesson:BN-C447-postano}



\begin{quote}
Before the new one: say the Bengali for remorse, then the Bengali for repentant.
\end{quote}

\subsection*{You'll want to know: \bn{পস্তানো}}
\textbf{\bn{পস্তানো}} — \emph{pôstāno} — \textquotedblleft{}to rue; to regret (something done)\textquotedblright{}.

\begin{quote}
\bn{সস্তা} \bn{জুতো} \bn{কিনে} \bn{পরে} \bn{পস্তালাম।}
\end{quote}

\begin{quote}
\emph{sôstā juto kine pôre pôstālām.}
\end{quote}

\begin{quote}
I bought cheap shoes and regretted it later.
\end{quote}

Past (he/she …): \bn{তিনি} \bn{পরে} \bn{পস্তালেন}, \emph{tini pôre pôstālen}.

\begin{cousinweb}
\textbf{\bn{পস্তানো}} (\emph{pôstāno}) is listed by CDIAL for Bengali under a reconstructed Sanskrit \emph{paścōttāpa}, \textquotedblleft{}burning afterwards\textquotedblright{}, that is, remorse (CDIAL 8010).
\end{cousinweb}

\begin{grammarlens}[title={What you've built}]
One more word for describing what you lived and felt.
\end{grammarlens}

\subsection*{Your turn}
\begin{itemize}
\raggedright
  \item \emph{Say it:} \emph{pôstāno}
  \item \emph{Say it:} \emph{pôstāno}, once more
  \item \emph{Say it:} \emph{sôstā juto kine pôre pôstālām.} — I bought cheap shoes and regretted it later.
  \item \emph{Recall:} say the Bengali for remorse, then the Bengali for repentant, then say \emph{pôstāno} again
\end{itemize}

\begin{tcolorbox}[breakable,colback=teal!4,colframe=teal!35!black,title={Before you move on}]
What does \bn{পস্তানো} mean? (\textquotedblleft{}To rue; to regret (something done)\textquotedblright{}.) Say it once more.
\end{tcolorbox}
\section[\texorpdfstring{\bn{স্বীকার} \bn{করা} (svīkār kôrā)}{স্বীকার করা (svīkār kôrā)}]{\bn{স্বীকার} \bn{করা} (svīkār kôrā) — to admit; to acknowledge}

\label{lesson:BN-C447-svikar-kora}



\begin{quote}
Before the new one: say the Bengali for repentant, then the Bengali for to rue.
\end{quote}

\subsection*{You'll want to know: \bn{স্বীকার} \bn{করা}}
\textbf{\bn{স্বীকার} \bn{করা}} — \emph{svīkār kôrā} — \textquotedblleft{}to admit; to acknowledge\textquotedblright{}.

\begin{quote}
\bn{অরুণ} \bn{ভুলটা} \bn{স্বীকার} \bn{করল।}
\end{quote}

\begin{quote}
\emph{ôrun bhulṭā svīkār kôrlô.}
\end{quote}

\begin{quote}
Arun admitted the mistake.
\end{quote}

Past (he/she …): \bn{তিনি} \bn{সত্যিটা} \bn{স্বীকার} \bn{করলেন}, \emph{tini sôttiṭā svīkār kôrlen}.

\begin{cousinweb}
\textbf{\bn{স্বীকার} \bn{করা}} (\emph{svīkār kôrā}) is Sanskrit \emph{svīkāra}, \textquotedblleft{}making one's own, accepting\textquotedblright{} (Monier-Williams), with \emph{kôrā}, \textquotedblleft{}to do\textquotedblright{} (CDIAL 2814).
\end{cousinweb}

\begin{grammarlens}[title={What you've built}]
One more word for describing what you lived and felt.
\end{grammarlens}

\subsection*{Your turn}
\begin{itemize}
\raggedright
  \item \emph{Say it:} \emph{svīkār kôrā}
  \item \emph{Say it:} \emph{svīkār kôrā}, once more
  \item \emph{Say it:} \emph{ôrun bhulṭā svīkār kôrlô.} — Arun admitted the mistake.
  \item \emph{Recall:} say the Bengali for repentant, then the Bengali for to rue, then say \emph{svīkār kôrā} again
\end{itemize}

\begin{tcolorbox}[breakable,colback=teal!4,colframe=teal!35!black,title={Before you move on}]
What does \bn{স্বীকার} \bn{করা} mean? (\textquotedblleft{}To admit; to acknowledge\textquotedblright{}.) Say it once more.
\end{tcolorbox}
